% !PRINTED 93
\weberpage{93}
%<<CONT
oder, da die ganzen Zahlen $x_k$ beliebig sind, die Summen
\[
\sum_{1,\mu}^{h}\frac{\alpha_h\,\gamma_{k,h}}{c_h}\qquad k=1,2,\ldots \nu
\]
ganze Zahlen sein müssen. Setzt man dann nach (15) und (11) \S.\,21
\[
\gamma_{k,h}=c_{k,h}\,e_{k,h},\quad c_h=\delta_{k,h}\,c_{k,h},\quad i_k=\delta_{k,h}\,i_{k,h},
\]
so folgt, dass
\[
\sum^{h}\frac{\alpha_h\,e_{k,h}}{\delta_{k,h}}=\sum^{h}\frac{\alpha_h\,i_{k,h}\,e_{k,h}}{i_k}
\]
ganze Zahlen sein müssen, und diese Forderung ist gleichbedeutend mit dem System der Congruenzen
\[
\sum^h \alpha_h\,i_{k,h}\,e_{k,h}\equiv 0\ (\mathrm{mod}\ i_k),
\]
oder, ausführlicher geschrieben
\weq{3}{\begin{aligned}
&e_{1,1}i_{1,1}\alpha_1+e_{1,2}i_{1,2}\alpha_2+\cdots+e_{1,\mu}i_{1,\mu}\alpha_\mu\equiv 0\ (\mathrm{mod}\ i_1)\\
&e_{2,1}i_{2,1}\alpha_1+e_{2,2}i_{2,2}\alpha_2+\cdots+e_{2,\mu}i_{2,\mu}\alpha_\mu\equiv 0\ (\mathrm{mod}\ i_2)\\
&\cdots\cdots\cdots\cdots\cdots\cdots\cdots\cdots\cdots\cdots\cdots\cdots\cdots\cdots\cdots\cdots\cdots\cdots\cdots\cdots\\
&e_{\nu,1}i_{\nu,1}\alpha_1+e_{\nu,2}i_{\nu,2}\alpha_2+\cdots+e_{\nu,\mu}i_{\nu,\mu}\alpha_\mu\equiv 0\ (\mathrm{mod}\ i_\mu). % [sic: Weber prints i_\mu; the row pattern requires i_\nu.]
\end{aligned}}
Durch diese Congruenzen sind, wie es sein muss, die Zahlen $\alpha_1, \alpha_2, \ldots \alpha_\mu$ nur nach den Moduln $c_1, c_2, \ldots c_\mu$ bestimmt, weil $i_{k,h}c_h=i_k c_{k,h}$ durch $i_k$ theilbar ist, und \gesperrt{um die Indices aller Zahlen der Gruppe $\mathfrak{A}$ zu finden, hat man einfach alle Lösungen der Congruenzen (3) aufzusuchen.}

\begin{center}
\rule{4cm}{0.4pt}
\end{center}
